% A1 — draft v0.1 (18/07/2026). Target: Information Processing & Management
% (alt.: ESWA; EMNLP industry). Class: article (converter p/ elsarticle na submissão).
\documentclass[11pt]{article}
\usepackage[T1]{fontenc}
\usepackage[utf8]{inputenc}
\usepackage[english]{babel}
\usepackage{booktabs}
\usepackage{amsmath}
\usepackage[round]{natbib}
\usepackage[margin=2.8cm]{geometry}
\usepackage{microtype}
\usepackage{xcolor}
\newcommand{\todo}[1]{\textcolor{red}{[TODO: #1]}}

\title{Large language models as labeling oracles for 621-category product
classification in Portuguese:\\ an instrumented evaluation of quality and cost}

\author{Gilsiley Henrique Darú$^{1}$ \and Gustavo Valentim Loch$^{1}$\\[2pt]
\normalsize $^{1}$Graduate Program in Numerical Methods in Engineering (PPGMNE),\\
\normalsize Federal University of Paraná (UFPR), Curitiba, Brazil}
\date{Draft v0.1 --- \today}

\begin{document}
\maketitle

\begin{abstract}
Large language models (LLMs) are increasingly used as labeling oracles in
place of human annotators, but the evidence base concentrates on tasks with
tens to a few hundred classes and rarely treats cost or the measurement
instrument itself as objects of study. We evaluate six LLM oracles on
zero-shot classification of short retail product descriptions in Portuguese
into a closed space of 621 categories --- a granularity three to ten times
beyond prior product-classification studies --- under a fully instrumented
protocol: structured output constrained to the label space, batch annotation
with paired calibration, prefix-cache cost accounting, provenance recording
of the (model, provider, date) triple, and statistics designed for paired
comparisons (Wilson intervals, exact McNemar tests). Four top oracles form a
plateau at 77--83\% accuracy, within $\approx 7$ points of a fully
supervised ceiling trained on 250{,}000 labels, at costs of
US\$~0.000--0.92 per thousand labels --- with the second-cheapest paid model
delivering 96\% of the best accuracy at 8.5\% of its cost, and a zero-cost
oracle statistically tied with it. Macro-F1 of the plateau oracles
\emph{exceeds} that of a light supervised baseline trained on the full
label budget, because LLMs are strongest precisely in rare classes. Error
anatomy shows nearly half of oracle errors reflect catalog conventions
rather than product misunderstanding. Two findings concern the instrument
itself: without output constraints, 6.8\% of responses become spurious
phrasing errors that a naive protocol counts against the model; and the
same model served by two providers disagrees significantly with itself
($p<0.001$) on identical instances, making serving provenance part of any
credible measurement. A prompt ablation with explicit boundary rules
improves the production distribution ($+4.6$ points, $p=0.012$) while
severely degrading the stratified tail ($-10.8$ points, $p<0.001$),
a cautionary double-edged result for prompt optimization. All code, prompts
and per-instance artifacts are public.
\end{abstract}

\noindent\textbf{Keywords:} large language models; data annotation; text
classification; short text; active learning; annotation cost; Portuguese.

\section{Introduction}
\label{sec:intro}

Supervised text classification remains bounded by the cost of labeled data.
The recent proposal to replace human annotators with large language models
promises orders-of-magnitude cost reduction \citep{Gilardi2023}, and a
rapidly growing literature explores LLMs as annotators inside active
learning loops \citep{Zhang2025,Rouzegar2024,Kholodna2024,Bayer2024}. Yet
three aspects of this evidence base limit its transfer to production
settings. First, \emph{granularity}: published product-classification
evaluations operate on tens to $\approx 370$ categories
\citep{Roumeliotis2025,Gholamian2024}; industrial catalogs routinely exceed
several hundred. Second, \emph{cost}: most studies report label counts, not
monetary cost, and none we are aware of instruments the pricing levers that
dominate real spending --- batch amortization of the prompt prefix and
provider-side prefix caching. Third, and least discussed, the
\emph{measurement instrument}: accuracy numbers are reported as properties
of a model, when they are in fact properties of a (model, prompt, output
contract, serving provider, date) configuration.

This paper evaluates six LLM oracles --- spanning commercial APIs, MaaS
platforms and a zero-cost evaluation tier --- on zero-shot classification of
Brazilian retail product descriptions (4--50 characters, uppercase,
aggressive abbreviations) into a closed space of \textbf{621 categories}.
The evaluation is \emph{instrumented} end to end: every annotation records
its instrument mode, batch size, token counts and serving provenance; all
comparisons are paired on identical instances; and the gold standard's own
error rate is audited and propagated as a measurement ceiling.

We ask four research questions. \textbf{RQ1}: what accuracy and Macro-F1 do
LLM oracles reach at this granularity, and are their differences
statistically real? \textbf{RQ2}: what does a thousand labels actually cost,
and how do batching and caching change it? \textbf{RQ3}: what is the
anatomy of oracle error? \textbf{RQ4}: how much of a measured ``error rate''
is an artifact of the measurement instrument itself?

Contributions: (i) the highest-granularity instrumented evaluation of
LLM-as-annotator for product classification to date, in a lower-resource
language; (ii) a reusable measurement protocol (output contract, paired
batch calibration, cache-aware cost accounting, provenance triple) with
public artifacts; (iii) two instrument-level findings --- the output-format
artifact and serving-provider divergence --- that affect any future oracle
evaluation; and (iv) an actionable error anatomy plus a prompt-ablation
result showing that boundary rules derived from aggregate error profiles
transfer asymmetrically between the production and tail distributions.

\section{Related work}
\label{sec:related}

\paragraph{LLMs as annotators.} \citet{Gilardi2023} established that
zero-shot ChatGPT can outperform crowdworkers at a fraction of the cost;
subsequent work integrated LLM annotation into active learning loops
\citep{Zhang2025}, routed hard cases to humans by confidence
\citep{Rouzegar2024}, and demonstrated batch annotation economics in
low-resource languages \citep{Kholodna2024}. This line inherits the older
noisy-annotator tradition \citep{Snow2008,Sheng2008}. Evaluations at
product-catalog scale are rarer: \citet{Roumeliotis2025} compare zero-shot
LLMs on 248 categories; \citet{Gholamian2024} test robustness on up to 370
leaf categories. Our 621-way closed space extends this frontier while
adding paired statistics and cost instrumentation absent from both.

\paragraph{Annotation noise.} Oracle errors become label noise for any
downstream learner; taxonomies and robustness results
\citep{Frenay2014,Song2023NoisyLabels} distinguish uniform from structured
noise, and recent benchmarks show \emph{real} annotation noise to be more
harmful than synthetic corruption \citep{Merdjanovska2024NoiseBench,
Raczkowska2024AlleNoise}. Our error anatomy (RQ3) characterizes where LLM
oracle noise falls in this spectrum.

\paragraph{Portuguese retail classification.} \citet{Machado2026RetailPt}
classify $\approx$100k Portuguese supermarket titles into ECOICOP
categories with fine-tuned BERT variants, relying on an ad-hoc
human-labeling pipeline; they do not evaluate LLM oracles. Community
surveys report deployment complexity and cost uncertainty as the main
barriers to LLM-assisted annotation \citep{Romberg2025Reassessing} ---
precisely the gap an instrumented cost protocol addresses.

\section{Task, data and gold standard}
\label{sec:data}

The task is single-label classification of short product descriptions from
Brazilian fiscal receipts (e.g., \texttt{CERV BRAHMA LT 350ML}) into a
closed catalog of 621 product-type categories. The dataset comprises
$\approx$250k labeled descriptions \citep{Daru2024Dissertacao}; texts are
4--50 characters, uppercase, with domain-specific abbreviations. The label
space exhibits severe long-tail imbalance, which motivates Macro-F1 as the
primary quality metric alongside accuracy.

\paragraph{Gold-standard audit.} Since the reference labels were produced
by a human process, we audited label conflicts (identical normalized text
with divergent labels) and re-scored all oracles (a) excluding conflicting
instances and (b) under multi-gold scoring. Accuracy shifts by at most
$+0.7$ points and no ranking changes; the known gold error bounds the
measurement ceiling at $\approx 99.3\%$ and is reported alongside all
results rather than ignored.

\paragraph{Evaluation samples.} All oracles annotate \emph{the same}
instances (fully paired design). \textbf{S-rand} ($n=1{,}000$, simple
random) preserves the production distribution and serves RQ1/RQ2/RQ4; its
size caps the Wilson 95\% half-width at $\approx 3$ points (2.5 in the
observed plateau). \textbf{S-strat} (3 per class, $n=1{,}863$) provides
per-class support for Macro-F1 and confusion analysis (RQ3): with 621
classes, a random sample leaves most classes unrepresented, and three
instances per class is the minimum that separates systematic (0/3, 3/3)
from sporadic (1/3, 2/3) class error.

\section{Measurement instrumentation}
\label{sec:instrument}

\paragraph{Output contract.} Three instrument modes are defined and
recorded per annotation: \textbf{enum} (production mode) --- structured
output with the label constrained at decoding time to the 621-value schema
(provider-enforced structured outputs); \textbf{json-prompt} --- for
providers without structured output support, the category list and JSON
format are instructed in the prompt, with strict post-validation and
explicit accounting of invalid labels; \textbf{free} --- structured output
\emph{without} label constraint, a controlled replica of the naive
instrument used solely to quantify its artifact (RQ4).

\paragraph{Prompt and batching.} The prompt contextualizes the domain
(fiscal-receipt abbreviations), instructs type-level classification
ignoring brand and packaging, and discourages misuse of the sentinel
rare-class label. The static prefix (instructions + schema) precedes the
variable items, exploiting provider prefix caches. Annotation is batched
(numbered items in one call, indexed vector response); batch size is
calibrated in a paired pre-study (1 vs.\ 10 vs.\ 25 items on identical
instances): exact McNemar detects no degradation ($b{=}1$ vs.\ $b{=}10$:
$p=0.58$), while batching cuts cost per label by up to $10\times$ through
prefix amortization --- consistent with the batch-annotation economics of
\citet{Kholodna2024}. We adopt $b=10$ (OpenAI) and $b=25$ (MaaS), recorded
per annotation.

\paragraph{Provenance and statistics.} Every annotation records the
(model, provider, mode, batch, temperature, date) tuple; temperature is 0
throughout. Accuracies carry Wilson 95\% intervals
\citep{Wilson1927,Brown2001}; model pairs are compared with McNemar's test
on discordant pairs (exact binomial when $b+c<25$)
\citep{McNemar1947,Dietterich1998}, the recommended regime for classifiers
evaluated once on a common sample.

\section{Results}
\label{sec:results}

\subsection{RQ1: quality at 621-way granularity}

Table~\ref{tab:main} reports the paired evaluation.
Four oracles form a plateau at 77--83\% accuracy on both samples.
Within-plateau differences are small but partly real: deepseek-v4-pro is
significantly better than v4-flash on S-strat (McNemar $b{=}43$, $c{=}16$,
$p<0.001$) and statistically tied with gpt-4o ($p=0.061$). The small
gpt-4o-mini collapses on the stratified sample (45.7\%): its errors
concentrate systematically in rare classes, which S-strat exposes by
construction. No oracle reaches the 85\% pre-registered adoption gate; the
best zero-shot result sits $\approx 7$ points below a fully supervised
ceiling trained on 250k labels (89.6\%, \citealp{Daru2024Dissertacao}).
Conversely, plateau Macro-F1 (0.78--0.80 on S-strat) \emph{exceeds} the
0.70 of a light supervised model trained on the full label budget: world
knowledge compensates exactly where supervision is scarcest --- the long
tail.

\begin{table}[htb]
\centering
\caption{Oracle performance on paired samples (temperature 0, calibrated
batching). Enum mode for OpenAI; json-prompt for MaaS/NIM. CI: Wilson 95\%.}
\label{tab:main}
\small
\begin{tabular}{llrrrr}
\toprule
\textbf{Oracle} & \textbf{Sample} & \textbf{Acc.} & \textbf{95\% CI} &
\textbf{Macro-F1} & \textbf{Invalid} \\
\midrule
deepseek-v4-pro   & S-rand  & 82.1\% & [79.6, 84.4] & 0.785 & 0.1\% \\
gpt-4o            & S-rand  & 80.9\% & [78.4, 83.2] & 0.788 & 0.0\% \\
deepseek-v4-flash & S-rand  & 78.3\% & [75.6, 80.8] & 0.750 & 0.7\% \\
nemotron-3-ultra (free) & S-rand & 77.9\% & [75.2, 80.4] & 0.752 & 1.2\% \\
glm-5.2           & S-rand  & 77.3\% & [74.6, 79.8] & 0.742 & 0.0\% \\
gpt-4o-mini       & S-rand  & 61.3\% & [58.2, 64.3] & 0.518 & 0.0\% \\
\midrule
deepseek-v4-pro   & S-strat & 82.6\% & [80.8, 84.2] & 0.793 & 0.0\% \\
gpt-4o            & S-strat & 82.1\% & [80.3, 83.8] & 0.797 & 0.0\% \\
deepseek-v4-flash & S-strat & 81.6\% & [79.7, 83.3] & 0.781 & 0.9\% \\
glm-5.2           & S-strat & 80.9\% & [79.1, 82.6] & 0.782 & 0.0\% \\
nemotron-3-ultra (free) & S-strat & 79.6\% & [77.7, 81.4] & 0.770 & 2.6\% \\
gpt-4o-mini       & S-strat & 45.7\% & [43.4, 48.0] & 0.427 & 0.0\% \\
\bottomrule
\end{tabular}
\end{table}

\subsection{RQ2: cost, cache and the serving provider}

\begin{table}[htb]
\centering
\caption{Cost per thousand labels (US\$, prices at execution, Jul/2026)
and prefix-cache hit rate on S-rand.}
\label{tab:cost}
\small
\begin{tabular}{lrrr}
\toprule
\textbf{Oracle} & \textbf{US\$/1k} & \textbf{Cache} & \textbf{Acc.} \\
\midrule
nemotron-3-ultra (free) & 0.000 & ---  & 77.9\% \\
deepseek-v4-flash & 0.035 & ---  & 78.3\% \\
gpt-4o-mini       & 0.051 & 91\% & 61.3\% \\
glm-5.2           & 0.249 & ---  & 77.3\% \\
deepseek-v4-pro   & 0.410 & ---  & 82.1\% \\
gpt-4o            & 0.923 & 88\% & 80.9\% \\
\bottomrule
\end{tabular}
\end{table}

Costs within the same accuracy plateau span more than an order of
magnitude (Table~\ref{tab:cost}): deepseek-v4-flash delivers 78--82\% at
US\$~0.035/1k --- $26\times$ cheaper than gpt-4o at comparable quality ---
and two to three orders of magnitude below typical human annotation costs
\citep{Gilardi2023}. Prefix caching reaches 88--95\% hit rates where
supported, halving effective input-token cost. Cost, however, is not the
only operational constraint: provider throughput limits (3 requests/min on
the MaaS tier; 18/min on a free aggregator tier) shape wall-clock time as
much as price --- a cheap, slow oracle can be dominated by an expensive,
parallel one when the deadline binds.

The zero-monetary-cost regime proved competitive: nemotron-3-ultra served
via NVIDIA's evaluation API ties deepseek-v4-flash on S-rand (McNemar,
$p=0.76$). Two operational caveats are findings in themselves. (i) Free is
not operationally free: a first attempt through an aggregator with a
50-requests/day quota projected six days of wall time for the same
workload. (ii) \emph{The serving provider is part of the instrument}: the
same nemotron model served by the aggregator disagrees significantly with
itself served by NVIDIA on identical instances (paired subset $n=850$,
$p<0.001$). Oracle measurements should therefore record not only the
model, but who serves it and when --- each number is a photograph of a
(model, provider, date) triple, not a property of model weights.

\subsection{RQ3: anatomy of oracle error}

Among the 324 errors of the best oracle (v4-pro, S-strat): \textbf{31\%}
are confusions between \emph{sibling categories} sharing a head word
(\textit{antisseptico bucal} $\to$ \textit{enxaguante bucal});
\textbf{17\%} have an umbrella class (``other \ldots'') as gold while the
model picks the specific class, and 7\% the reverse; $\approx 2\%$ involve
the sentinel class or invalid output; the remainder are genuine semantic
errors. Nearly half of oracle error therefore reflects not product
misunderstanding but ignorance of \emph{this catalog's conventions} ---
boundaries between neighbors and aggregation policy, learnable only from
labeled examples. Two consequences: it motivates the prompt ablation of
Section~\ref{sec:prompt}, and it predicts oracle noise is
\emph{structured} (concentrated in adjacent pairs), the more benign regime
for downstream learners \citep{Frenay2014,Song2023NoisyLabels}.

\subsection{RQ4: the instrument effect}

Holding model (gpt-4o-mini), instances (S-rand) and prompt constant, and
switching only the output contract --- enum vs.\ free --- yields accuracy
61.3\% vs.\ 60.5\%, but invalid-label rates of \textbf{0.0\% vs.\ 6.8\%}.
The instrument's effect appears less in point accuracy than in the
\emph{accounting of error}: without output constraints, nearly 7\% of
responses are out-of-space phrasings that a naive protocol counts as model
error, deflating measurements and contaminating comparisons between models
with different paraphrasing propensities. Enum-constrained decoding with
strict validation eliminates this artifact class by construction.

\subsection{Prompt as an instrument variable}
\label{sec:prompt}

Since $\approx 48\%$ of RQ3 errors are convention errors, we test whether
prompt engineering can recover them in the cheapest cached model
(gpt-4o-mini): variant \textbf{v4a} adds ten explicit boundary rules
derived from the aggregate error profile; \textbf{v4b} adds ten
\emph{invented} boundary examples (never sampled from evaluation data);
both are paired against the base prompt on the first 500 items of each
sample. On S-rand, rules recover part of the convention error: $+3.8$
points (v4a, $p=0.045$) and $+4.6$ points (v4b, $p=0.012$). On S-strat the
same rules \emph{destroy} performance: $-7.0$ and $-10.8$ points
($p<0.001$), with invented examples significantly worse than rules alone
($p=0.0013$). Rules distilled from aggregate error profiles push the model
toward the conventions of \emph{frequent} classes --- exactly the behavior
the stratified sample punishes. Prompt optimization for labeling oracles
must therefore be validated on both the production and the tail
distribution before adoption; a gain measured only on the natural
distribution can mask severe tail regression.

\section{Discussion}
\label{sec:discussion}

\paragraph{Choosing an oracle is a portfolio decision.} At 621-way
granularity no oracle passes a 85\% adoption gate, yet the plateau's
Macro-F1 beats a fully supervised light baseline. The practical reading
for label-scarce pipelines: the LLM oracle is not a replacement for
supervision but a complementary specialist for the long tail, at
US\$~0.00--0.41 per thousand labels. The 26$\times$ price spread inside a
statistical tie makes oracle selection a measurable engineering decision,
not a brand choice --- and our protocol makes it auditable.

\paragraph{Measure the instrument, not just the model.} Two of our
findings would silently bias any naive replication: the 6.8\%
output-format artifact and the significant self-disagreement across
serving providers. Both argue that credible LLM-annotation papers should
publish their output contract, serving provenance and per-instance
artifacts --- as we do --- and that leaderboard-style comparisons without
this metadata are fragile.

\paragraph{Threats to validity.} Single dataset and language pair
(mitigated by public artifacts and a replication design on AlleNoise
\citep{Raczkowska2024AlleNoise} as future work); API prices dated at
execution (ratios are more stable than absolute values); boundary rules in
the ablation derive from this catalog's error profile (the method
generalizes; the rules do not); gold-standard noise bounds all accuracies
at $\approx 99.3\%$, audited and propagated.

\section{Conclusion}
\label{sec:conclusion}

We provide the highest-granularity instrumented evaluation of LLMs as
labeling oracles to date: 621 closed categories, Portuguese short text,
six oracles across three cost regimes, fully paired statistics, and
cache-aware cost accounting. LLM oracles plateau at 77--83\%, seven points
under a 250k-label supervised ceiling, while beating it on Macro-F1; a
zero-cost oracle ties the best cost-efficiency tier; nearly half of oracle
error is catalog convention, recoverable by prompt rules only at the
tail's expense; and the measurement instrument itself --- output contract
and serving provider --- moves results by margins that exceed several
model-to-model differences. The protocol and artifacts are public and
directly reusable for oracle selection in any closed-taxonomy annotation
pipeline.

\paragraph{Artifacts.} Code, prompts, per-instance annotations and
analysis notebooks: \texttt{github.com/GHDaru/activelearning}
(\texttt{experiments/e0}, \texttt{e0p}).
\todo{DOI/Zenodo release + decisão sobre publicação integral da base}

\bibliographystyle{plainnat}
\bibliography{refs}

\end{document}
